\usepackage{fontspec}
\usepackage{unicode-math}
%% s35 (Gate-9 G05).  These were the family names STIX Two Text / STIX Two Math / Menlo, resolved
%% through the operating system's font database; none of the three is in the TeX tree, so a clean
%% Linux build of the archive stopped.  The same faces are now taken from the TeX distribution
%% (stix2-otf) and named by FILE, so kpathsea resolves them and the operating system is not asked.
\setmainfont{STIXTwoText-Regular.otf}[
  BoldFont       = STIXTwoText-Bold.otf,
  ItalicFont     = STIXTwoText-Italic.otf,
  BoldItalicFont = STIXTwoText-BoldItalic.otf]
\setmathfont{STIXTwoMath-Regular.otf}
%% Latin Modern Mono was tried first and missed 69 glyphs over twelve characters, all inside
%% \texttt{}: the subscripts and Greek letters of Lean identifiers.  DejaVu Sans Mono ships with
%% TeX Live and covers all twelve, checked in the font's cmap before this line was written.
\setmonofont{DejaVuSansMono.ttf}[
  BoldFont   = DejaVuSansMono-Bold.ttf,
  ItalicFont = DejaVuSansMono-Oblique.ttf,
  Scale      = 0.80]
\usepackage[margin=1in]{geometry}
\usepackage{lineno}
\linenumbers
\modulolinenumbers[1]
\usepackage{longtable,booktabs,array}
\usepackage{amsmath}
\usepackage[hidelinks]{hyperref}
% s37 (Gate-11 N18).  The shipped PDF carried no document metadata at all: a desk editor's
% reader showed an empty Title and Author.  These are set from the same strings as the
% typeset front matter; if the title changes, this block changes with it.
\hypersetup{
  pdftitle={Structure and computability of the single-quadratic existential fragment over the integers},
  pdfauthor={Hiroki Fukui},
  pdfsubject={Mathematical logic; Diophantine definability; existential fragments over the integers},
  pdfkeywords={Diophantine definability; existential fragments; quadratic equations; finite fibres; definable functions; computability},
  pdfcreator={XeLaTeX with hyperref}}
\setlength{\parskip}{0.4\baselineskip}
\setlength{\parindent}{0pt}
\sloppy
\raggedbottom
\emergencystretch=3em

%% The 39 characters that occur in RUNNING TEXT and are absent from the text
%% font. Each is made active and routed to the math font. The list was
%% generated from the "Missing character" lines of a previous run, not guessed;
%% the build gate below requires that list to come back empty.
\makeatletter
\newcommand{\champmap}[2]{\catcode`#1=\active \begingroup\lccode`~=`#1
  \lowercase{\endgroup\def~}{\ensuremath{#2}}}
\makeatother
\champmap{∈}{\in}
\champmap{∉}{\notin}
\champmap{≤}{\leq}
\champmap{≥}{\geq}
\champmap{≠}{\neq}
\champmap{≡}{\equiv}
\champmap{∣}{\mid}
\champmap{∤}{\nmid}
\champmap{→}{\to}
\champmap{↦}{\mapsto}
\champmap{⇒}{\Rightarrow}
\champmap{⇐}{\Leftarrow}
\champmap{⟹}{\Longrightarrow}
\champmap{⟺}{\Longleftrightarrow}
\champmap{∎}{\blacksquare}
\champmap{∃}{\exists}
\champmap{∀}{\forall}
\champmap{⊆}{\subseteq}
\champmap{⊊}{\subsetneq}
\champmap{∞}{\infty}
\champmap{√}{\surd}
\champmap{∧}{\wedge}
\champmap{∨}{\vee}
\champmap{∩}{\cap}
\champmap{⊕}{\oplus}
\champmap{∅}{\emptyset}
\champmap{∘}{\circ}
\champmap{ℓ}{\ell}
\champmap{⟨}{\langle}
\champmap{⟩}{\rangle}
\champmap{𝒫}{\mathcal{P}}
\champmap{ᵀ}{^{\mathsf{T}}}
\champmap{ᵏ}{^{k}}
\champmap{ᵐ}{^{m}}
\champmap{ᵢ}{_{i}}
\champmap{ⱼ}{_{j}}
\champmap{ₖ}{_{k}}
\champmap{⋯}{\cdots}
\champmap{ᵖ}{^{p}}      % added s31: the glyph gate caught this one in the proof of Theorem 3.4   % added s31: the glyph gate caught this one on a new Appendix F proof

%% ---- native-source additions (s36, Gate-10 N01-N09) -------------------------------------
%% One shared counter per section for every numbered statement kind, so that Theorem 2.1,
%% Theorem 2.2, Remark 2.3, ... come out exactly as the Markdown numbered them; the converter
%% recomputed every number and would have refused a mismatch.
\usepackage{amsthm}
\theoremstyle{plain}
\newtheorem{thm}{Theorem}[section]
\newtheorem{theorem}[thm]{Theorem}
\newtheorem{lemma}[thm]{Lemma}
\newtheorem{proposition}[thm]{Proposition}
\newtheorem{corollary}[thm]{Corollary}
\newtheorem{conjecture}[thm]{Conjecture}
\theoremstyle{definition}
\newtheorem{definition}[thm]{Definition}
\newtheorem{example}[thm]{Examples}
\newtheorem{notation}[thm]{Notation}
\newtheorem{question}[thm]{Question}
\theoremstyle{remark}
\newtheorem{remark}[thm]{Remark}
%% s38 (Gate-12 T01): cross-references are \cref against the kind:number labels; names are
%% capitalised so that the printed text equals the hard-coded text it replaced.
\usepackage[capitalise,noabbrev,nameinlink]{cleveref}
\crefname{thm}{Theorem}{Theorems}
\crefname{theorem}{Theorem}{Theorems}
\crefname{lemma}{Lemma}{Lemmas}
\crefname{proposition}{Proposition}{Propositions}
\crefname{corollary}{Corollary}{Corollaries}
\crefname{definition}{Definition}{Definitions}
\crefname{example}{Example}{Examples}
\crefname{remark}{Remark}{Remarks}
\crefname{section}{Section}{Sections}
\crefname{appendix}{Appendix}{Appendices}
%% unnumbered forms (Appendix F Parts I-IV state their results without numbers)
\theoremstyle{plain}
\newtheorem*{theorem*}{Theorem}
\newtheorem*{lemma*}{Lemma}
\newtheorem*{proposition*}{Proposition}
\newtheorem*{corollary*}{Corollary}
\theoremstyle{definition}
\newtheorem*{definition*}{Definition}
\newtheorem*{example*}{Examples}
\theoremstyle{remark}
\newtheorem*{remark*}{Remark}
%% Appendix F Part V numbers its results V.1 ... V.6 on one counter
\newcounter{Part}
\theoremstyle{plain}
\newtheorem{theoremPart}[Part]{Theorem}
\newtheorem{lemmaPart}[Part]{Lemma}
\newtheorem{propositionPart}[Part]{Proposition}
\newtheorem{corollaryPart}[Part]{Corollary}
\theoremstyle{definition}
\newtheorem{definitionPart}[Part]{Definition}
\theoremstyle{remark}
\newtheorem{remarkPart}[Part]{Remark}
%% displayed Unicode formulas (the Markdown's four-space-indented lines)
\newenvironment{displaytext}{\par\medskip\begin{quote}\raggedright\setlength{\parskip}{0pt}}{\end{quote}\medskip}
